\SourceLeaf{139}
Erhöhung der Witwenpensionen wurde beschlossen, ein Gutachten von \textsc{Kritter} einzuholen\footnote{Es scheint nicht, dass etwas darüber festgesetzt wäre, \textit{in welchem Maasse} \textsc{Kritter}’s Rath in Anspruch genommen werden solle. \textsc{Kaestner} war in der Versammlung nicht gegenwärtig.}. Doch erklärte man sich schon \textit{gegen} die im Rescripte beantragte grössere Pension für die sechs ältesten Witwen, als welche schon durch das V.’sche Legat bevorzugt seien.

Die Consultation des p. \textsc{Kritter}, welche durch ein im Concept bei den Acten befindliches Schreiben P.’s geschah, war in der That nur eine sehr beschränkte. Nach einer bloss in ganz allgemeinen Umrissen gehaltenen Uebersicht der Haupteinrichtungen der Witwenkasse werden \textsc{Kritter} nur zwei Fragen vorgelegt: I) Nach welchem Grundsatz und Verhältniss in einer derartigen Gesellschaft das Maximum der Witwen bestimmt werden müsse und II) um wieviel dies Maximum der zu verabreichenden Witwenpensionen noch vergrössert werden müsse, wenn im Falle des Nichtvorhandenseins einer Witwe, oder nach dem frühern Ableben derselben die Pensionsberechtigung auf etwa vorhandene Waisen übergehe, und fortdaure, bis das jüngste Kind das Alter von 20 Jahren erreicht habe. Die dermalige Anzahl aller Interessenten der Kasse, und der darunter befindlichen Verehelichten, wird gerade eben so wie in der oben angeführten Anlage des Rescripts, zu 36 und 26 angegeben, und zugleich bemerkt, dass diese Zahlen und ihr Verhältniss veränderlich und schwerlich einer Wahrscheinlichkeitsregel zu unterwerfen seien; zum Schluss folgt das Anerbieten, dass wenn der Befragte noch einige weitere Data aus den bisherigen Erfahrungen über das Verhältniss der Participanten und der Witwen nöthig haben sollte, solche sogleich mitgetheilt werden würden.

\textsc{Kritter}’s Antwort, oder sein ‘Gutachten’, vom 19. Julius, lege ich in einer vollständigen Abschrift bei\footnote{Gutachten über einige mir vorgelegte Fragen, die Göttingische Universitäts-Witwenkasse betreffend.

Ad I. Das Collegium der Herrn Professoren in Göttingen hat schon seit der Errichtung der Universität über 50 Jahre lang existirt, so dass man schon vor 10 Jahren die höchste Zahl der Witwen haben konnte, welche nach den Gesetzen der Sterblichkeit auf etwa 22 oder 24 verehelichte Professoren vorhanden sein mussten, nemlich 12 Witwen. Dieses hat sich auch laut der mir mitgetheilten Liste von dem Anwachs der jährlich vermehrten Zahl der Witwen gezeigt und würde sich noch besser gezeigt haben, wenn das Collegium der Herrn Professoren aus 100 Personen hätte bestehen können. Da aber nach der Natur der Sache diese Anzahl nur um den \(4^{\text{ten}}\) Theil von 100 stark gewesen, so war es auch natürlich, dass die Zahl der Witwen vom Jahre 1775 bis 1780 mehr als 12 und vom Jahre 1781 bis 1792 weniger als 12 betragen, weil bei kleinen Zahlen die Ordnung der Sterblichkeit nicht so wie bei grossen Zahlen eintreten kann. Indessen muss man dennoch annehmen, dass die höchste Zahl der Witwen in einem Durchschnitt von etwa 10 Jahren beständig etwa halb so gross sein werde, als die Zahl der verehelichten Herrn Professoren.

Ad II. Da es nunmehr gewünscht wird, dass die hinterlassenen Waisen-Familien einer gestorbenen Witwe oder auch eines gestorbenen Witwers bis zum vollendeten \(20^{\text{sten}}\) Jahre eine gleiche Pension wie eine Witwe bekommen möchten, so kann ich nur aus der Erfahrung bei der Bremischen Witwenkasse etwas davon bestimmen. Die Bremische Gesellschaft hatte etwa den \(6^{\text{ten}}\) Theil der Witwenpensionen mehr zu erwarten, da sie die Waisen-Familien gleich einer Witwe zu pensioniren und bis zum vollendeten \(18^{\text{ten}}\) Jahre des jüngsten Kindes damit fortzufahren beschloss. Da aber die Waisen-Familien der Herrn Professoren bis zum vollendeten \(20^{\text{sten}}\) Jahre des jüngsten Kindes dieses geniessen sollen, so müssen doch wohl gegen 12 Witwenpensionen über 2 Waisenpensionen gerechnet werden.

Ich halte es also für rathsam, dass man zu 12 als dem Maximo der Witwenzahl noch wenigstens 2 addire, so dass der Divisor, worin die jährlichen Einkünfte von den Fonds der Casse und sonst dividirt werden, auf 14 gesetzt werde. Da nun jetzt nur 10 Witwen vorhanden sind, so darf man die sämtlichen jährlichen Revenüen der Casse nicht unter diese 10 vertheilen, weil sonst die in der Folge hinzukommenden Witwen würden verkürzt werden, sondern man muss in 14 dividiren, so wird innerhalb 10 Jahren der Fonds auf etwa 25 Witwenportionen verstärkt sein, wovon die Zinsen noch auf eine Pension mehr als vorhin zureichen werden, so dass man alsdann, wenn das wahre Maximum der Witwen- und Waisenpensionen eintreten wird, 15 Pensionen wird bezahlen können; und sollte auch etwas übrig bleiben, so könnte vorzüglich armen Witwen etwas zugelegt werden.

Auf mögliche Unglücksfälle bei den belegten Capitalien, Verlust an Zinsen und andern Ausgaben musste doch auch wohl etwas gerechnet werden.

Göttingen den 19. Juli 1794.\hfill J. A. \textsc{Kritter}.}.
